\documentclass{article}

\title{Proof of Soundness Preservation under Place Substitution in Workflow Nets}
\author{} % Intentionally blank
\date{}   % Intentionally blank

\begin{document}
% \maketitle % Uncomment if a title is desired

\section*{Solution}
\textbf{Note: } The "soundness" from the notes refers to 1-soundness, but the one in the paperr instead referes to $N$-soundness for each $N$ in the first case the theorem is incorrect, so I assume that we're talking about the $N$-soundness for each $N$. \\
Let $N_1=(P_1,T_1,F_1,i_1,o_1)$ and $N_2=(P_2,T_2,F_2,i_2,o_2)$ be sound WF-nets. Assume $P_1 \cap P_2 = \emptyset$ and $T_1 \cap T_2 = \emptyset$ (disjoint sets of nodes). Let $N'=(P',T',F',i',o')$ be the WF-net constructed by substituting place $p \in P_1$ with the entirety of $N_2$. 
% Specifically, the set of places $P'=(P_1 \setminus \{p\}) \cup P_2$, and the set of transitions $T'=T_1 \cup T_2$. Arcs in $F_1$ originally incident to $p$ are reconnected: for any $x$ such that $(x,p) \in F_1$, we have $(x,i_2) \in F'$; for any $y$ such that $(p,y) \in F_1$, we have $(o_2,y) \in F'$. If $p \notin \{i_1,o_1\}$, then the start place $i'$ of $N'$ is $i_1$, and the end place $o'$ of $N'$ is $o_1$. (The argument holds analogously if $p = i_1$ or $p = o_1$, where $i'$ or $o'$ would be $i_2$ or $o_2$ respectively).

Check that the net satisfies the definition of soundness: 
\begin{itemize}
\item For each token put in pin, one and only one token eventually appears in pout,
\item When the token appears in pout, all other places are empty, and
\item For each transition (i.e, task), it is possible to move from the initial state to a state in which this transition is enabled.

The third condition is obvious from the fact that both nets are sound: for any marking, the tokens inside the subnet can be moved to the end since it's sound. And if they Are at the "out" of the subnet, we can obtain a sequence of transitions that moves them to the "out".

\end{document}
